\begin{scanpage}{0103}{86}
\[\left.\begin{matrix}f(3)\\g(3)\end{matrix}\right\}=\frac{\sin\frac\pi n}n\left(\frac1{\cos\frac\pi n-\cos\frac{\vartheta-\vartheta_0}n}\mp\frac1{\cos\frac\pi n-\cos\frac{\pi+\vartheta+\vartheta_0}n}\right).\tag{15.13}\]
Expressions (15.10) satisfy the reciprocity principle --- that is, they do not change their values if one interchanges $\vartheta$ and $\vartheta_0$. When $\vartheta=\vartheta_0$, they change into the previous Expressions (15.02).

Equations (15.02) and (15.10) describe the radiation from the currents flowing only on part of the cylinder's surface: on the one end (when $z=\frac l2$) and on half of the lateral surface $(-\pi\leq\psi\leq0)$. Moreover, these expressions do not take into account the nonuniform part of the current caused by the curvature of the cylindrical surface. Therefore, they must be refined with values of $\vartheta$ and $\vartheta_0$ which are close to $\frac\pi2$ and $\pi$. However, in the case $\vartheta=\vartheta_0$ --- that is, in the direction towards the source --- these corrections may be neglected if the parameters $ka$ and $kl$ are sufficiently large. Numerical calculations performed by us on the basis of Equations (15.02) show that this evidently may be done already when $ka=\pi$ and $kl=10\pi$. The graphs of the functions $\frac{\sigma_E}{\pi a^2}=|\overline\Sigma(\vartheta)|^2$ and $\frac{\sigma_H}{\pi a^2}=|\Sigma(\vartheta)|^2$ constructed for this case in Figures 25 and 26 agree with the experimental curve\textsuperscript{(1)} (the dashed line): the position of the maxima and minima basically agree, and the number of diffraction fringes is the same. For the purpose of illustrating the effect of the ends, we constructed a graph of the effective scattering area for those same values of $ka$ and $kl$, taking into account only the uniform part of the current on the cylindrical surface (Figure 27). A comparison of Figures 25, 26 and 27 shows that the effect of the ends begins to appear when $\vartheta=120^\circ$.

\vfill\noindent\rule{.5\linewidth}{.4pt}

Footnote (1) appears on page 89.
\end{scanpage}
